\section{总结与延伸}

\subsection{讲者的核心总结}

Zane Durante 在课程结尾总结了三个要点：
\begin{enumerate}
    \item Vanilla RNN 概念简洁但性能有限
    \item LSTM 等变体通过门控机制选择性传递信息，缓解了梯度消失问题
    \item RNN 的核心思想（循环状态、线性复杂度）正在通过状态空间模型回归研究前沿
\end{enumerate}

\subsection{全课知识图谱}

\begin{center}
\begin{tikzpicture}[node distance=0.8cm, every node/.style={draw, rounded corners, minimum width=3.5cm, minimum height=0.7cm, font=\small}]
\node (seq) {序列建模问题};
\node (rnn) [below=of seq] {Vanilla RNN};
\node (bptt) [below=of rnn] {BPTT 与截断 BPTT};
\node (grad) [below=of bptt] {梯度消失/爆炸};
\node (lstm) [below=of grad] {LSTM / GRU};
\node (ssm) [below=of lstm] {状态空间模型};
\draw[->, thick] (seq) -- (rnn);
\draw[->, thick] (rnn) -- (bptt);
\draw[->, thick] (bptt) -- (grad);
\draw[->, thick] (grad) -- (lstm);
\draw[->, thick] (lstm) -- (ssm);
\node[draw=none, right=0.8cm of seq, text width=5cm, font=\scriptsize] {one-to-one/many, many-to-one/many};
\node[draw=none, right=0.8cm of rnn, text width=5cm, font=\scriptsize] {$h_t = \tanh(W_{hh}h_{t-1} + W_{xh}x_t)$};
\node[draw=none, right=0.8cm of bptt, text width=5cm, font=\scriptsize] {共享权重、梯度累加、分块训练};
\node[draw=none, right=0.8cm of grad, text width=5cm, font=\scriptsize] {$\tanh'$ 衰减 + $W_{hh}$ 奇异值问题};
\node[draw=none, right=0.8cm of lstm, text width=5cm, font=\scriptsize] {Cell State Highway, 门控机制};
\node[draw=none, right=0.8cm of ssm, text width=5cm, font=\scriptsize] {Mamba, RWKV: 线性复杂度};
\end{tikzpicture}
\end{center}

\subsection{关键 Takeaways}

\begin{importantbox}{五条核心原则}
\begin{enumerate}
    \item \textbf{隐藏状态是 RNN 的灵魂}：它在时间步之间传递信息，编码序列的``记忆''
    \item \textbf{参数共享使模型泛化}：所有时间步共享权重，使 RNN 能处理任意长度的输入
    \item \textbf{梯度消失是长序列的天敌}：$\tanh$ 和权重矩阵的反复累乘导致远距离梯度信号衰减
    \item \textbf{LSTM 的 cell state 是解药}：无激活函数的直通路径，类似 ResNet 的跳跃连接
    \item \textbf{RNN 并未过时}：状态空间模型正在以线性复杂度挑战 Transformer 的二次方开销
\end{enumerate}
\end{importantbox}

\end{document}
